% !TeX root = ../quiver-avatar-moduli.tex
% 2026-11-22 / v0.5（OU z2座標とPicard根の明示対応、Kummer障害を追加）
\section{既存GITと内在的モジュライの比較}
\label{sec:qa-git-comparison}

本節は Okawa--Ueda \cite[Theorem 1.1]{OU2014} の
quintuple と conifold potential の GIT 商を、
本稿の有限 foundation と内在的 $\HHstack$ に対する比較問題へ接続する。
二つの既存 GIT 商自体を本稿の新しい成果とはしない。
また、表示を忘れた $\HHstack$ への射は
定義~\ref{def:qa-compatible-chart} の適合する領域でのみ構成し、
GIT compactification 全体で自動的に存在するとは主張しない。

\subsection{既存の次数四射と三つの忘却}

$V_0,V_1,V_2,V_3$ を二次元、$W$ を一次元の $k$-ベクトル空間とし、
\[
 T=(V_0\otimes V_1\otimes V_2\otimes V_3)\otimes W^\vee,\qquad
 \widetilde G=\prod_{i=0}^3\GL(V_i)\times\GL(W)
\]
と置く。Okawa--Uedaの生のquintuple表示スタックを
$\mathfrak M_{\mathrm{qui}}^{\mathrm{raw}}=[T/\widetilde G]$ と書く。
$T$ の零元を含む表示スタックの記法であり、
$\HHstack$ を定義するものではない。
同論文のpotential表示を $\mathfrak M_{\mathrm{pot}}$ とし、
その compact GIT 商を $\overline M_{\mathrm{pot}}$、
quintuple側の商を $\overline M_{\mathrm{qui}}$ とする。
既知結果は
\[
 \overline M_{\mathrm{pot}}\simeq\PP(1,2,3,4),\qquad
 \overline M_{\mathrm{qui}}\simeq\PP(2,4,4,6),\qquad
 \overline M_{\mathrm{pot}}\longrightarrow\overline M_{\mathrm{qui}}
 \text{ が有限次数 }4
\]
である。右端の射は潜在的なSerre平方根の忘却射ではなく、
\emph{potentialからquintupleへの表示の射}である。

本研究では区別すべき射が少なくとも三つある。
potentialからquintupleへの射、quintupleから四項foundationへの
表示上の射、そして適合foundationから pointed exact 対象
$(\mathcal C,\OO)$ への忘却射である。
最後の射では、quintupleの標識、ruling、Serre平方根とその同型、
ambientを保持しない。これらの射のどの組合せが
相対スタック射として存在するかは、対象と自然同型を個別に調べる必要がある。

\subsection{スカラー慣性の核}

\begin{proposition}[生のquintuple表示の四次元スカラー核]
\label{prop:qa-ou-scalar-gerbe}
次の中央部分群スキームは $\widetilde G$ に含まれる。
\[
 K=\left\{
 \bigl(\lambda_0\mathrm{id}_{V_0},\ldots,
       \lambda_3\mathrm{id}_{V_3},
       (\lambda_0\lambda_1\lambda_2\lambda_3)\mathrm{id}_{W}\bigr)
 :\lambda_i\in\Gm
 \right\}\simeq\Gm^4.
 \tag{GIT-K}\label{eq:qa-ou-scalar-kernel}
\]
$K$ は $T$ 全体に自明に作用する。従って
\[
 [T/\widetilde G]\longrightarrow[T/(\widetilde G/K)]
 \tag{GIT-rig}\label{eq:qa-ou-rigidification}
\]
は $K$ を相対慣性とする fppf gerbe である。

さらに、quintupleから対応する $\mathbb Z$-代数のtails圏と
添字零の構造対象を作る通常の構成が実際の相対族を与え、
その対象が $\HHstack$ の公理を満たす領域では、
$K$ の全元は pointed exact 自己同値の\emph{恒等射}へ写る。
従って、この領域の生の表示スタックから $\HHstack$ への
標準的な忘却射は表現可能ではない。
少なくともこのスカラー慣性を消さずに
全忠実な相対比較射を得ることはできない。
\end{proposition}
\begin{proof}
$\widetilde G$ のスカラー部分の作用は
$T$ 上で
\[
 (\lambda_0,\ldots,\lambda_3,\mu)\cdot t
   =(\lambda_0\lambda_1\lambda_2\lambda_3)\mu^{-1}t
\]
である。従って式~\eqref{eq:qa-ou-scalar-kernel} は
$T$ への作用の核に含まれる中央な四次元トーラスである。
$\widetilde G/K$ は群スキームとして存在し、
$\widetilde G\to\widetilde G/K$ は忠実平坦なので、
$K$ の作用が自明である場合の商スタックの一般的な構成から
式~\eqref{eq:qa-ou-rigidification} はgerbeになる。

後半は表現を変えるスカラーが圏に与える作用を直接調べる。
quintupleを使って構成した $\mathbb Z$-代数を $D$ とし、
その隣接Hom空間を $V_i$（または採用する双対規約）で表す。
スカラー変更は、各頂点の単元を変えない
\[
 D_{ij}\longrightarrow D_{ij},\qquad
 a\longmapsto\gamma_i^{-1}\gamma_j\,a
\]
という次数ごとの共役に帰着する。
ここで隣接空間での所定のスカラーに合うよう
$\gamma_i\in k^\times$ を漸化的に選ぶ。
四つの隣接スカラーの積が一次関係線 $W$ のスカラーと一致する
ことが、quintupleの関係を保つ条件である。
右graded $D$-加群の次数 $i$ の成分に
$\gamma_i^{-1}$ を掛ける自然変換は、この共役による
twist functorと恒等 functorの自然同型を与える。
したがってtailsへの降下でも恒等と自然同型であり、
添字零の構造対象の同型も同じスカラーで規格化できる。
相対族でも単元を用いて同じ計算ができる。

一方、元の表示スタックでは任意の $K$ の元は
各点の自己同型として残る。
忘却による相対慣性の核が $K$ を含むため、
射が表現可能であるための必要条件（相対慣性の自明性）に反する。
\end{proof}

\begin{remark}[GITの次数と慣性を混同しない]
\label{rem:qa-ou-degree-four-not-root}
Okawa--Ueda \cite[Theorem 1.1(4)]{OU2014} の次数四は、
compact GIT 商の有限射の一般ファイバーの次数である。
式~\eqref{eq:qa-ou-scalar-kernel} の群は
quintuple表示の自己同型に既に存在し、
pointed忘却では消える。
したがってその次数四を、標識を忘れた
$\HHstack$ の一般ファイバーの位数、
あるいは可換点でのSerre平方根の個数と同一視できない。
適合する族の範囲、表示側の剛体化、全pointed同型の像と核、
および閉軌道の選択が別途必要である。
\end{remark}

\subsection{相対ファイバーの群oidとしての公式}

\begin{proposition}[忘却射の幾何ファイバーと慣性核]
\label{prop:qa-forgetful-homotopy-fiber}
$F:\mathfrak X\to\mathfrak Y$ を $k$-上の groupoid stack の射とし、
$y\in\mathfrak Y(k)$ を取る。
幾何ファイバーの groupoid
$\mathfrak X_y=\mathfrak X\times_{\mathfrak Y}\operatorname{Spec}k$
の対象は、$x\in\mathfrak X(k)$ と同型
$\alpha:F(x)\xrightarrow{\sim}y$ の組である。
その点 $(x,\alpha)$ の自己同型群は
\[
 \operatorname{Aut}_{\mathfrak X_y}(x,\alpha)
 =\ker\bigl(\operatorname{Aut}_{\mathfrak X}(x)
       \longrightarrow\operatorname{Aut}_{\mathfrak Y}(F(x))\bigr).
 \tag{Fib-I}\label{eq:qa-fiber-inertia}
\]
さらに、$F(x)\simeq y$ となる $\mathfrak X(k)$ の同型類を
一つずつ選ぶと、$\mathfrak X_y(k)$ の同型類の集合は
\[
 \coprod_{[x]:\,F(x)\simeq y}
 \operatorname{Aut}_{\mathfrak Y}(y)\Big/
  \alpha\,F(\operatorname{Aut}_{\mathfrak X}(x))\,\alpha^{-1}
 \tag{Fib-C}\label{eq:qa-fiber-cosets}
\]
と同一視される。右辺は群の商ではなく、一般には剰余\emph{集合}である。
従って同型類のファイバーの濃度、スタックの相対慣性、
GITの有限射の次数はそれぞれ異なる量である。
\end{proposition}
\begin{proof}
二つの対象 $(x,\alpha),(x',\alpha')$ の間の射は
$f:x\to x'$ で $\alpha'F(f)=\alpha$ を満たすものである。
$x=x'$、$\alpha=\alpha'$ とすれば
式~\eqref{eq:qa-fiber-inertia} が直ちに従う。
固定した $x$ について同型 $\alpha:F(x)\to y$ を一つ選ぶと、
他の全ての同型は $g\alpha$、
$g\in\operatorname{Aut}_{\mathfrak Y}(y)$ の形である。
二つの選択は $\operatorname{Aut}_{\mathfrak X}(x)$ の像による作用で
同型になるため、式~\eqref{eq:qa-fiber-cosets} が得られる。
\end{proof}


\subsection{次数四射の $S_4$-monodromy と Picard 平方根}

\begin{proposition}[OU次数四射の一般Galois閉包]
\label{prop:qa-ou-fourfold-s4-monodromy}
標数零の代数閉体上で、Okawa--Ueda \cite[\S5, (5.8)]{OU2014}
のGIT有限射を、対角化による稠密な一般開で考える。
同論文の表示は
\[
 \PP^3/S_3\longrightarrow
 \PP^3/\bigl(S_3\ltimes(\mathbb Z/2)^3\bigr)
 \tag{OU-4}\label{eq:qa-ou-fourfold-diagonal}
\]
であり、$\PP^3$ では同時符号反転が自明に作用する。
従って有効な符号変更群と全群はそれぞれ
\[
 \Gamma=(\mathbb F_2)^3/
   \langle(1,1,1)\rangle\simeq(\mathbb F_2)^2,\qquad
 G=\Gamma\rtimes S_3\simeq\operatorname{AGL}_2(\mathbb F_2)
 \simeq S_4.
\]
$S_3$ は $\Gamma$ の三つの非零元を置換する標準的な
$\operatorname{GL}_2(\mathbb F_2)$-作用を持つ。
一般の射~\eqref{eq:qa-ou-fourfold-diagonal} は次数四で、
\emph{Galois射ではなく}、その関数体拡大のGalois閉包の群は $S_4$ である。
$\Gamma$ を自明化する表示の方へ引き戻した場合にのみ、
$\Gamma$-torsorとしての四点の符号変更を直接見ることができる。
\end{proposition}
\begin{proof}
$\mathbb F_2^3$ の標準基底の剰余類を $e_1,e_2,e_3$ とする。
商 $\Gamma$ では $e_1+e_2+e_3=0$ なので、
三つはまさに非零元である。$S_3$ がこれらを置換する作用は忠実で、
$\operatorname{Aut}_{\mathbb F_2}(\Gamma)
=\operatorname{GL}_2(\mathbb F_2)$ は位数六なので同型となる。
$\Gamma$ の四元上への平行移動とこの線型作用を合わせた
$G$ は四元上の忠実で推移的な位数 $4\cdot6=24$ の置換群、
従って $S_4$ である。
$H=S_3$ は原点の安定化群で、$[G:H]=4$。
標数零の一般対角座標はこの有限群作用に対し自明な
幾何安定化群を持つので、
関数体拡大 $k(\PP^3)^H/k(\PP^3)^G$ の次数は四である。
$H$ の $G$ における core（全共役の共通部分）は、
四元上の推移作用の核と一致し、忠実性から自明である。
従ってこの非正規拡大のGalois閉包は $k(\PP^3)$ であり、
Galois群は $G\simeq S_4$ となる。
最終のtorsorに関する文は、
$H$ の選択を消した一般の基底では
$\Gamma$ が一定群として作用する四重Galois被覆を
無条件に得られない、という意味である。
置換群の全24元とcoreは
\path{computations/verify_ou_fourfold.py} で再現できる。
\end{proof}

\begin{proposition}[適合Picard平方根の幾何fiberの濃度]
\label{prop:qa-picard-roots-degeneration-count}
論文Bの非線型geometric cubic quintupleの相対曲線
$C\to S$ と $N\in\Pic^0_{C/S}(S)$ を考え、
同論文 \cite[\texttt{lem:picard-halves}]{QuadricModuli2026}
の平方根空間
\[
 T=S\times_{\Pic^0_{C/S},[-2]}\Pic^0_{C/S}
 \tag{PR-4}\label{eq:qa-picard-root-torsor}
\]
を取る。これは $S$ 上 étale 全射である。
滑らかな楕円曲線の locus $S_{\mathrm{sm}}$ 上では
$T_{S_{\mathrm{sm}}}\to S_{\mathrm{sm}}$ は
$\Pic^0_{C/S}[2]$-torsorであり、有限 étale 次数四である。
各代数閉幾何点での平方根の個数は、曲線が
滑らかな楕円曲線、既約一節点の種数一曲線、
既約一尖点の種数一曲線である場合、それぞれ
\[
 \begin{array}{c|ccc}
 C_{\bar s}&\text{smooth elliptic}&\text{nodal}&\text{cuspidal}\\ \hline
 \Pic^0(C_{\bar s})&E&\mathbb G_m&\mathbb G_a\\
 \#T_{\bar s}&4&2&1
 \end{array}
 \tag{PR-count}\label{eq:qa-picard-root-count}
\]
となる。従って滑らかなfiberと節点fiberが同じ連結基底族に
現れるとき、$T\to S$ はその族上で有限 étale
\emph{次数四のまま}延長することはできない。
\end{proposition}
\begin{proof}
$\Pic^0_{C/S}$ は滑らかな相対一次元可換群代数空間で、
$[-2]$ は étale 全射であることを論文Bの
\cite[\texttt{lem:picard-halves}]{QuadricModuli2026} が示す。
滑らかな種数一のfiberのPicard零成分は楕円曲線の
相対Jacobianであり、$[-2]$ は有限 étale
次数 $2^{2}=4$ のisogenyである。
そのpullbackは四点の$E[2]$-torsorとなる。

一節点曲線の正規化は $\PP^1$ で、節点の二つの原像の
fiberの貼合せスカラーにより
$\Pic^0\simeq\mathbb G_m$ となる。
$[-2]$ は $z\mapsto z^{-2}$ なので、
代数閉体の各非零値の逆像は二点である。
一尖点曲線では正規化との貼合せの一次jetの加法群から
$\Pic^0\simeq\mathbb G_a$ となる。
$[-2]$ は $z\mapsto-2z$、従って一意の平方根を持つ。
いずれも標数零では étale である。
最後の主張は有限 étale 射の次数が連結基底上で
局所定数となることから従う。
\end{proof}

\begin{remark}[四点の抽象的な一致と比較すべきtorsor class]
\label{rem:qa-ou-picard-fourfold-not-identical}
OU側の $\Gamma\simeq\mathbb F_2^2$ と
楕円曲線上の $E[2](\bar k)\simeq\mathbb F_2^2$ は、
$S_3=\operatorname{GL}_2(\mathbb F_2)$ の
三非零元に対する作用も含めて同じ抽象モデルを持つ。
これは次数四の一致を説明する\emph{候補}であって、
OUの有限射とPicard平方根空間の同型を証明するものではない。
特にOU側は一般に$S_4$-monodromyの非Galois四重射であり、
Picard側は変動する群スキーム$E[2]$のtorsorである。
二つが同じ基底上に引き戻され、
$E[2]$の局所系とOUの符号変更局所系が同定できた後に、
両torsorの相対同型を確認する必要がある。
標識、自己同型群、さらにpointed同値による商は
その後に初めて比較できる。
節点・尖点における式~\eqref{eq:qa-picard-root-count} は
Picard根空間そのもののfiberの個数であり、
OU射の分岐fiberや$\HHstack$への忘却fiberを
数えたものではない。
\end{remark}

\begin{proposition}[OU楕円四つ組の $z_2$ と適合Picard根]
\label{prop:qa-ou-elliptic-half-coordinate}
滑らかな楕円曲線 $E$ に原点 $o$ を取り、
Okawa--Ueda \cite[\S4, (4.12)--(4.13)]{OU2014}
の楕円四つ組を
\[
 (E,\mathcal L_0,\mathcal L_1,\mathcal L_2)
 =
 (E,\OO_E(2o),\OO_E(2z_1),\OO_E(2z_2))
\]
と書く。論文Bの
$N=\mathcal L_2\mathcal L_0^{-1}\in\Pic^0(E)$ に対する
$a^{-2}=N$ の解集合は
\[
 \left\{
 \OO_E(o-z_2)\otimes\xi:\xi\in E[2]
 \right\}
 \tag{OU-PR}\label{eq:qa-ou-elliptic-half-coordinate}
\]
である。従ってOUの楕円四つ組の birational 表示において、
$\mathcal L_0,\mathcal L_1,\mathcal L_2$ を固定したまま
$z_2$ に2-torsionを加える操作は、
\emph{対応するPicard平方根の四つの選択}に一致する。
これは式~(4.13)の $(\mathbb Z/2)^4$ のうち
$z_2$ に作用する一つの $(\mathbb Z/2)^2$ の因子である。

ただし、式~\eqref{eq:qa-ou-fourfold-diagonal} の
potential GITからquintuple GITへの次数四の
\emph{四つのsheet}がこの $z_2$ 因子のtorsorと同型であることは、
この命題の結論ではない。
\end{proposition}
\begin{proof}
$N=\OO_E(2z_2-2o)$ であるから
$a_0=\OO_E(o-z_2)$ は
$a_0^{\otimes(-2)}=N$ を満たす。
二つの根$a,a_0$の比は二乗が自明な
$\Pic^0(E)[2]=E[2]$ の元であり、逆にその全元は解を与える。
$z_2$ の取り換え $z_2\mapsto z_2+t$（$2t=0$）は
$\OO_E(2z_2)$ を変えず、$a_0$ を
$\OO_E(-t)$ 倍する。従って四つの選択と自由かつ推移的に対応する。
OUの式~(4.11)--(4.13)は $z_1,z_2$ の独立した
2-torsion平行移動を含むので、最後の因子の同定も従う。
\end{proof}

\begin{lemma}[Picard平方根のKummer class]
\label{lem:qa-picard-kummer-torsor}
滑らかな楕円曲線の相対族 $E\to S$ と
$N\in E(S)$ に対し、平方根空間
$T_N=S\times_{E,[-2]}E$ の
$E[2]$-torsor class は、étale siteのKummer列
\[
 0\longrightarrow E[2]\longrightarrow E
 \xrightarrow{[-2]}E\longrightarrow0
\]
から得られる連結写像による
$\delta_{[-2]}(N)\in H^1_{\mathrm{\acute et}}(S,E[2])$
である。
従って、同じ基底上でOUの四重射に
$E[2]$-torsor構造を与える局所系の同型を
別途構成できた場合に限り、そのtorsorとPicard根空間が同型となる
必要十分条件は二つの $H^1$ class の一致である。
\end{lemma}
\begin{proof}
$[-2]$ は滑らかな楕円群スキーム上の有限étale isogenyで
核は$E[2]$なので、示した短完全列はétale層として完全。
$N$ の局所持上げの差から得るČech 1-cocycleこそ
$T_N$ のtorsor classである。
比較先にも同じ構造群スキームを同定した後は、
torsorの同型類の分類が$H^1$で与えられることから従う。
\end{proof}

\subsection{一般非線型 cubic の実際の相対比較射}

\begin{theorem}[geometric quintuple から pointed 相対スタックへの射]
\label{thm:qa-geometric-cubic-relative-forgetting}
$\mathscr G^{\mathrm{nl}}$ を、式~\eqref{eq:qa-ou-scalar-kernel}
のテンソル表示空間 $T\setminus\{0\}$ のうち
Van den Bergh の\emph{非線型 geometric quintuple}
を与える $\widetilde G$-不変開部分とする。
このとき局所有限型 $k$-スキーム上の fppf site において
\[
 \mathcal F_{\mathrm{geo}}:
 [\mathscr G^{\mathrm{nl}}/\widetilde G]
 \longrightarrow\HHstack,\qquad
 (V_0,V_1,V_2,V_3,W\hookrightarrow
       V_0\otimes\cdots\otimes V_3)
 \longmapsto
 (\operatorname{qgr}C,\pi e_0C)
 \tag{GIT-H2}\label{eq:qa-geometric-cubic-comparison-map}
\]
という\emph{実際の相対群oid射}が定義される。
ここで $D$ はquintupleのcubic $\mathbb Z$-代数、
右加群・向きの規約を論文Bに合わせて
$C_{ij}=D_{-j,-i}$ とする。
$\operatorname{qgr}C$ の相対heartは
module-object 基底変更・fppf descent を伴って構成される。

この射は $K\simeq\Gm^4$ の慣性を自明にする自然な
二同型を持ち、その剛体化を通じて
\[
 \overline{\mathcal F}_{\mathrm{geo}}:
 [\mathscr G^{\mathrm{nl}}/(\widetilde G/K)]
 \longrightarrow\HHstack
 \tag{GIT-H2-rig}\label{eq:qa-geometric-cubic-rigid-map}
\]
と因数分解する。

さらに、論文Bの相対Picard二等分空間
$T_{\mathrm{root}}\to\mathscr G^{\mathrm{nl}}$ で選んだ
二周期 $\beta$ と Serre 平方同型を使って得る
$\HHstack$-対象は、選択にかかわらず
式~\eqref{eq:qa-geometric-cubic-comparison-map}
の像の引戻しである。
特に滑らかな楕円曲線上の四つのPicard根は、
\emph{固定quintupleの四つの異なるpointed対象ではない}。

この定理は $\overline{\mathcal F}_{\mathrm{geo}}$ の
表現可能性・全忠実性・幾何ファイバーの有限性、
およびGIT compactification全体への延長を主張しない。
\end{theorem}
\begin{proof}
$\mathscr G^{\mathrm{nl}}$ 上の普遍quintupleは
相対的なcubic $\mathbb Z$-代数 $D$ を定める。
双対・添字反転により $C_{ij}=D_{-j,-i}$ とするのは
論文B \cite[\texttt{thm:geometric-cubic-relative-effectivity}]{QuadricModuli2026}
の右tails規約である。
同論文 \cite[\texttt{prop:general-cubic-canonical-tails}]{QuadricModuli2026}
により、二周期を選ぶ\emph{前}から
$\operatorname{qgr}C$ と $\pi e_0C$ は
全係数基底変更に整合する相対Abel族であり、
compact Morita理論・Serre双加群も基底上に構成される。

一方、同論文
\cite[\texttt{lem:picard-halves},
\texttt{thm:relative-compatible-root}]{QuadricModuli2026}
は普遍基底を被覆するétale局所で
実際の二周期 $\beta$ とそのSerre平方比較を構成し、
\cite[\texttt{thm:geometric-cubic-relative-effectivity}]{QuadricModuli2026}
が同じ相対圏と構造対象に対して
AZ性・全軌道AS4被覆を含む $\HHstack$ の条件を証明する。
これらは\emph{存在条件}で、追加標識ではなく
fppf局所で判定される。
従って基底上にもともとある
$(\operatorname{qgr}C,\pi e_0C)$ が
$\HHstack$ の実際の対象を定める。
ここでは選んだ $\beta$ を降下させたのではなく、
既に存在する圏と構造対象に対して、
必要な局所witnessの存在を確認している。
同じ議論は任意の局所有限型の非被約基底変更で成立する。

quintupleの同型は $\mathbb Z$-代数の同型を誘導し、
右tailsの相対heartおよび $\pi e_0C$ の同型を保つ。
この操作は恒等、合成、任意の基底変更と両立する。
従って $\widetilde G$-equivariantな族として
式~\eqref{eq:qa-geometric-cubic-comparison-map} が
fppf descentにより構成される。

命題~\ref{prop:qa-ou-scalar-gerbe} の頂点ごとの
共役による自然同型は、添字零でスカラーを1に規格化すれば
pointedである。各 $\gamma_i$ は四つのスカラーの
積による指標として選べるので、
$K$ の積に関する二同型の整合性も成立する。
従って $K$ の作用を剛体化した表示スタックへ
この射は降下する。
最後に根空間上で$\beta$を替えても
$\operatorname{qgr}C$と$\pi e_0C$は元の基底から
引き戻した\emph{同じ}対象であり、結論が従う。
\end{proof}

\begin{corollary}[quintuple の剛体化後の相対次元]
\label{cor:qa-quintuple-rigid-dimension}
前定理の剛体化された表示スタックは
$k$ 上smoothで、純次元三である。
その $\HHstack$ への射~\eqref{eq:qa-geometric-cubic-rigid-map}
は同じ局所次元三のsmooth Artinスタック間の射である。
この次元の一致は相対射の étale 性や
pointed同値群の全忠実性を単独で意味しない。
\end{corollary}
\begin{proof}
$\mathscr G^{\mathrm{nl}}$ は16次元テンソル空間の
smoothな開部分である。
$\dim\widetilde G=4\cdot\dim\GL_2+\dim\GL_1=17$、
$\dim K=4$ なので
$\dim(\widetilde G/K)=13$ である。
従ってそのquotient stackの次元は $16-13=3$、
smooth性は群と表示開のsmooth性から従う。
一方、論文B
\cite[\texttt{thm:global-h2-smooth-finite-type}]{QuadricModuli2026}
が $\HHstack$ のsmooth純次元三を与える。
両者の間の相対射の微分・慣性核の計算は
この数値だけからは従わない。
\end{proof}

\begin{remark}[Cで次に必要な具体的計算]
\label{rem:qa-git-next-fiber}
実際の相対射は定理~\ref{thm:qa-geometric-cubic-relative-forgetting} で構成した。
次の段階では、式~\eqref{eq:qa-ou-fourfold-diagonal} の
$S_4$-monodromyと式~\eqref{eq:qa-picard-root-torsor} の
相対$E[2]$-torsorを同じsmooth elliptic quintuple基底へ引き戻し、
二つの四重射のモノドロミーとtorsor classを比較する。
そのうえで一般geometric cubic quintupleの
適合する有限型開において
$\operatorname{Aut}(\mathcal C,\OO)$ と
標識を保存する自己同型群の像を計算する。
次に上の二式を用いて、rulingと四窓を忘れる
pointed相対ファイバーを求める。
その後に、potentialからquintupleへの既存次数四射と照合する。
楕円型の構造対象の格子・特殊な自己同型、
noncubic境界への特殊化は別の段階とする。
$\HHstack$ 全体のgood moduliの存在、
標識適合領域の開性・全表示との相対同値は
現時点では未証明として保持する。
\end{remark}
